\subsection{Detection Challenges} \label{sec:Detection_Challenges}

Many factors, including stellar variability and instrumental noise, can affect the accuracy of light-curve analysis in transit photometry. These effects not only influence parameter estimation but can also lead to false-positive detections. A false positive occurs when a signal is incorrectly interpreted as a transit even though it was not caused by an exoplanet-star system. Several physical and instrumental effects can produce signals that mimic transit patterns \citep{Morton_2016}.

\subsubsection{Stellar Variability}

Stars exhibit time-variable activity across their surfaces. Sunspots reduce the local surface brightness while faculae lead to an increase in the surrounding region. The timescale and amplitude are dependent on factors such as effective temperature \citep{2012IAUS..285..364M}. When a transit passes over these active regions, this can cause small fluctuations in the transit profile \citep{2013A&A...556A..19O}. Active regions evolve and rotate across the stellar surface, introducing time-dependent variability. This rotational modulation can produce broad, quasi-periodic trends in the light curve that may obscure or imitate shallow transit signals. As a result, stellar variability can affect both the detection of transits and the estimation of planetary parameters derived from the light curve \citep{Bonomo_2009}.

\subsubsection{Instrumentation} 

Missions focused on transit observations must rely on precise differential photometry. The \textit{Kepler} mission \citep{2016RPPh...79c6901B}, must detect signals as small as $\approx80$ parts per million (ppm) which corresponds to an Earth-sized planet transiting a Sun-like star \citep{2010ApJ...713L..92C}. Achieving the necessary sensitivity to detect these signals is limited by several instrumental factors. Photon and detector noise as well as background flux contamination all affect the observations and must be accounted for \citep{2011ApJS..197....6G}. 

In addition to fundamental photon and detector noise, there is astrophysical time series variability which introduces photometric deviations. As a result, the detection threshold cannot be described only using Gaussian white noise \citep{2006MNRAS.373..231P}. Often referred to as "red noise", this can arise from events such as thermal variations or stellar variability and may make it difficult to discern from genuine light curve trend. Transit depths and durations can be distorted by this noise, which can mimic shallow transit-like signals during analysis. As a result, the effective signal-to-red noise ratio relative to white-noise statistics is reduced, limiting confidence in exoplanet detections \citep{2006MNRAS.373..799C}. This forms a fundamental limitation of transit detection and thus careful de-trending, calibration and statistical analysis of photometric data is required. Beyond noise-related effects, transit detection is also limited by the temporal sampling of the observations.

In wide-field transit surveys, nearby stars can introduce additional challenges through photometric contamination. The photometric aperture used to measure a target star can also collect flux from neighbouring sources, producing a blended light curve. This crowding can dilute the measured transit or eclipse depth, making the event appear shallower than it truly is \citep{Morton_2016}. This is especially important for eclipsing binaries located within or near the photometric aperture, since their diluted eclipse signals can have depths comparable to planetary transits. These blended eclipsing binaries are therefore a major source of astrophysical false positives and are difficult to distinguish using photometry alone \citep{2011ApJ...738..170M}.

\subsubsection{Eclipsing Binaries}

A significant fraction of catalogued stars are found in binary star systems \citep{2010A&ARv..18...67T}. When the orbital plane is viewed close to edge-on, the two stars periodically pass in front of one another, producing eclipses in the observed light curve. The deeper primary eclipse occurs when the brighter star is occulted, while the secondary eclipse occurs when the fainter star is obscured. Systems with this geometry are known as eclipsing binaries (EBs).

Eclipsing binaries are a major source of false positives in transit surveys because their eclipse signatures can resemble planetary transits \citep{Morton_2016}. Depending on the geometry of the system, grazing eclipses can produce sharp, V-shaped dips, while more complete occultations can produce flatter minima \citep{2010exop.book...55W}. Grazing planetary transits can also appear V-shaped, making the distinction between planets and EBs difficult from photometry alone. In addition, if an EB is blended with light from a neighbouring star, the eclipse depth can be diluted until it becomes comparable to a planetary transit \citep{Morton_2016}. This is especially relevant for shallow eclipsing binaries, which can closely mimic exoplanet-like signals.

The eclipse depths in binary systems depend on the relative sizes and surface brightnesses of the two stars. In systems with similar eclipse depths, the light curve may be mistaken for a single repeating transit signal, leading to an incorrect orbital period estimate \citep{fp_first_habitable_zone}. Previous transit surveys have found both undiluted and diluted eclipsing binaries to be a dominant source of astrophysical false positives \citep{2008A&A...488L..47M, 2009IAUS..253..382A, 2013A&A...553A..30T}. This motivates the binary and multiclass classification tasks used in this work, where eclipsing binaries are treated both as a general false-positive class and, in the multiclass case, separated into shallow and deep eclipsing binary examples.